\section{The path a number takes}

This section is the map. It follows one number from a cell in a CSV file to a coefficient in
a printed equation, naming the class that does each piece of work. The arithmetic of the
fitting itself comes after; what matters here is the order, and what each stage is allowed
to hand on.

\subsection{Two equations, and which is which}

There are two recovery equations in play and they are not the same equation. Confusing them
is the easiest mistake to make with this project, so they are set out before anything else.

The first is the one asked about:

\begin{equation}
\label{eq:real-units}
\text{recovery \%} \;=\; 33.0192 + 3.16999\,J - 14.4538\,R - 0.0954576\,J^{2}
+ 0.837907\,R^{2}
\end{equation}

Here $J$ is current density in $\mathrm{mA/cm^2}$ and $R$ is the feed ratio in
$\mathrm{mol/mol}$, both in the units the bench reads them in. It is ordinary least squares
over the nine GO-10 runs, with all five estimable terms forced in at once. \textbf{No part of
this project prints equation~\ref{eq:real-units}.} It is a fit made by hand from the nine
rows the analyzer produces, to answer the question ``what do the raw numbers say if you fit
everything''.

The second is what \texttt{uv run limodel analyze} actually prints:

\begin{equation}
\label{eq:coded}
\text{recovery \%} \;=\; 11.5776 - 15.627\,x_2 + 6.08016\,x_2^{2},
\qquad
x_2 = \frac{\ln R - 1.395778}{1.039721}
\end{equation}

Three differences, all of them deliberate and all of them in the code:

\begin{enumerate}[nosep]
  \item \textbf{Logged, coded units.} \texttt{DesignPoints} takes the natural log of each
        setting, then centres and scales it so the runs span $-1$ to $+1$. Equation~\ref{eq:coded}
        is in $x_2$, not $R$.
  \item \textbf{Fewer terms.} \texttt{FittedOutcome} does not fit all five. It grows the model
        one term at a time, keeping a term only if it lowers the leave-one-out error. It kept
        two, and current density never got in at all.
  \item \textbf{No $J$ at all.} Equation~\ref{eq:real-units} has $J$ and $J^2$ with
        coefficients; equation~\ref{eq:coded} has neither, because adding $x_1$ made the
        out-of-sample error worse.
\end{enumerate}

Wherever this report says ``the code's equation'' it means equation~\ref{eq:coded}, and
``the five-term fit'' means equation~\ref{eq:real-units}.

\subsection{The whole pipeline}

Figure~\ref{fig:pipeline} is the shape of the program. \texttt{Analyzer.analyze} is three
lines long and those three lines are the whole order: preprocess, calculate, fit.

Two things are worth noticing in the chart. First, the data changes kind twice. The CSV rows
become \texttt{Run} objects at the loader and stay \texttt{Run} objects through the whole of
preprocessing; they become \texttt{Result} objects at the calculator and are never readings
again. Nothing downstream of \texttt{Calculator} can see a tank weight.

Second, the \texttt{Journal} sits to the side and every stage writes into it. It is one
record, made empty, handed to the \texttt{Analyzer}, and handed on down. Nothing in it
prints as it is written; the \texttt{Analyze} command reads the whole of it at the end. That
is why the report can open with everything that was left out before it shows a single result.

\begin{figure}[htbp]
\centering
\begin{tikzpicture}[
  node distance=5mm,
  stage/.style={draw, rounded corners, align=center, font=\small, text width=6.4cm,
                inner sep=4pt, minimum height=7.5mm},
  payload/.style={draw, align=center, font=\small, text width=6.4cm,
                  inner sep=4pt, minimum height=7.5mm},
  record/.style={draw, rounded corners, align=center, font=\small, text width=3.2cm,
                 inner sep=5pt},
  flow/.style={-Stealth, thick},
  note/.style={-Stealth, dashed, semithick}
]

\node[payload] (csv)
  {four CSV files in \texttt{data/}\\[1pt]
   \footnotesize \texttt{runs}, \texttt{stack\_readings}, \texttt{tank\_readings},
   \texttt{samples}};

\node[stage, below=of csv] (loader)
  {\texttt{DataLoader} on \texttt{CsvSource}\\[1pt]
   \footnotesize one pandera schema per table; a blank cell becomes \texttt{None}};

\node[payload, below=of loader] (runs)
  {21 \texttt{Run} objects\\[1pt]
   \footnotesize each holding \texttt{TankReading}s (with their \texttt{Sample}),
   \texttt{StackReading}s};

\node[stage, below=of runs] (prep)
  {\texttt{Preprocessor}\\[1pt]
   \footnotesize four \texttt{Step}s, in order (figure~\ref{fig:steps})};

\node[payload, below=of prep] (kept)
  {9 \texttt{Run} objects kept, some values cleared};

\node[stage, below=of kept] (calc)
  {\texttt{Calculator}\\[1pt]
   \footnotesize weights, lab results and the power log become quantities};

\node[payload, below=of calc] (results)
  {9 \texttt{Result} objects\\[1pt]
   \footnotesize two settings, two outcomes, five diagnostics};

\node[stage, below=of results] (surr)
  {\texttt{Surrogate}\\[1pt]
   \footnotesize per membrane: \texttt{DesignPoints}, \texttt{FittedOutcome},
   \texttt{BestSetupUnderTheGates}};

\node[payload, below=of surr] (models)
  {\texttt{SurrogateModels}\\[1pt]
   \footnotesize fitted equations, the gated search, the membranes refused};

\node[stage, below=of models] (cmd)
  {\texttt{Analyze} command\\[1pt]
   \footnotesize prints the journal, the run table, then the fits};

\foreach \a/\b in {csv/loader, loader/runs, runs/prep, prep/kept, kept/calc,
                   calc/results, results/surr, surr/models, models/cmd}
  \draw[flow] (\a) -- (\b);

\node[record] at ($(prep.east)+(4.3,-0.6)$) (journal)
  {\texttt{Journal}\\[2pt]
   \footnotesize one \texttt{Entry} per thing left out: \texttt{DROP} for a whole run,
   \texttt{SKIP} for one value, each with its id and its reason};

\draw[note] (prep.east) -- (journal.west);
\draw[note] (calc.east) -- (journal.west);
\draw[note] (journal.south) |- (cmd.east)
  node[pos=0.28, right, font=\footnotesize, align=left] {read at\\the end};

\end{tikzpicture}
\caption{The pipeline, from four CSV files to a printed equation. Solid arrows carry the
data; the dashed arrows into the journal are writes, and the one dashed arrow out of it is
the command reading the whole record at the end. Rounded boxes do work, square boxes are
what they hand on. The order is
\texttt{Analyzer.analyze}.}
\label{fig:pipeline}
\end{figure}

\subsection{The four preprocessing steps}

\texttt{Preprocessor} holds the four steps as a tuple and runs each on what the one before
returned. Figure~\ref{fig:steps} shows them with the counts on the real data.

The order is a decision. Dropping by membrane comes second so that the reading checks never
spend effort on a run the modelling throws away anyway. Dropping runs the lab cannot answer
for comes \emph{last}, after the checks, so that a run on its way out is still looked over
and its bad readings still get written down: the journal is a record of the whole dataset,
not only of the part that survived.

Two steps drop whole runs and refuse to return nothing: if every run fails,
\texttt{DropExcludedRuns} and \texttt{DropRunsOnAnotherMembrane} both raise rather than hand
on an empty list, because modelling nothing in silence would hide a misspelled membrane
name.

\begin{figure}[htbp]
\centering
\begin{tikzpicture}[
  node distance=7mm,
  stepbox/.style={draw, rounded corners, align=left, font=\small, text width=10.4cm,
               inner sep=5pt},
  tally/.style={font=\small},
  flow/.style={-Stealth, thick}
]

\node[tally] (in) {21 \texttt{Run} objects loaded};

\node[stepbox, below=of in] (s1)
  {\textbf{1. \texttt{DropExcludedRuns}} \hfill \emph{drops 0}\\[1pt]
   \footnotesize removes any run not marked \texttt{include} at the bench};

\node[stepbox, below=of s1] (s2)
  {\textbf{2. \texttt{DropRunsOnAnotherMembrane}} \hfill \emph{drops 11}\\[1pt]
   \footnotesize removes any run whose \texttt{cem} is not in
   \texttt{MEMBRANES\_MODELLED} (here \texttt{GO-10})};

\node[stepbox, below=of s2] (s3)
  {\textbf{3. \texttt{SkipSuspectReadings}} \hfill \emph{drops 0, clears 10 values}\\[1pt]
   \footnotesize removes no run; sets to \texttt{None} each single value a check cannot
   believe (figure~\ref{fig:inspectors})};

\node[stepbox, below=of s3] (s4)
  {\textbf{4. \texttt{DropRunsMissingAnEndSample}} \hfill \emph{drops 1 (AW)}\\[1pt]
   \footnotesize removes any run without both metals, in both the concentrate and the
   dilute, at its first and last sampled minute};

\node[tally, below=of s4] (out) {9 \texttt{Run} objects handed to the \texttt{Calculator}};

\foreach \a/\b/\l in {in/s1/{21 in}, s1/s2/{21 in}, s2/s3/{10 in}, s3/s4/{10 in},
                      s4/out/{9 out}}
  \draw[flow] (\a) -- (\b) node[midway, right, font=\footnotesize] {\l};

\end{tikzpicture}
\caption{The four preprocessing steps and what each one removes, with the counts from
\texttt{uv run limodel analyze} on the current data. Step 3 is the only one that changes a
run without removing it.}
\label{fig:steps}
\end{figure}

\subsection{Inside one step}

Figure~\ref{fig:inspectors} opens step~3, the only step that reaches inside a run. It is
three levels deep and each level has one job. \texttt{SkipSuspectReadings} chooses the
inspectors. An inspector chooses which readings it looks at and which checks they get: six
check classes, configured into 21 checks per run. A check decides, and decides nothing else:
it is handed the readings and answers with the one worst offender, the field to clear, and a
sentence saying why.

The clearing happens in exactly one place, \texttt{Inspector.\_clear\_what\_it\_rejects},
and in exactly this order: write to the journal, then set the attribute to \texttt{None}.
Nothing else in the codebase clears a reading. That loop also asks each check again after
every value it clears, because a median is a better median once a bad value has left it, and
it stops after a third of a series has gone: past half, a median starts following the bad
values and the check would begin clearing the readings that were right.

\begin{figure}[htbp]
\centering
\begin{tikzpicture}[
  scale=1, transform shape,
  insp/.style={draw, rounded corners, align=center, font=\footnotesize, text width=2.9cm,
              inner sep=4pt},
  wide/.style={draw, rounded corners, align=center, font=\small, text width=7.5cm,
               inner sep=5pt},
  leaf/.style={draw, rounded corners, align=left, font=\footnotesize, text width=6.1cm,
               inner sep=4pt},
  flow/.style={-Stealth, semithick}
]

\node[wide] (step) at (0,0)
  {\texttt{SkipSuspectReadings}\\[1pt]
   \footnotesize every run $\times$ every inspector};

\node[insp] (conc) at (-5.55,-1.9) {\texttt{ConcTank}\\\texttt{Inspector}\\[1pt] tank
  \texttt{conc}};
\node[insp] (dil)  at (-2.35,-1.9) {\texttt{DiluteTank}\\\texttt{Inspector}\\[1pt] tank
  \texttt{dil}};
\node[insp] (rin)  at ( 0.85,-1.9) {\texttt{RinseTank}\\\texttt{Inspector}\\[1pt] tank
  \texttt{rinse}};
\node[insp] (pow)  at ( 4.55,-1.9) {\texttt{PowerLog}\\\texttt{Inspector}\\[1pt] the
  \texttt{StackReading}s};

\foreach \b in {conc, dil, rin, pow} \draw[flow] (step) -- (\b);

\node[leaf] (tankchecks) at (-3.6,-4.7)
  {six checks, built once in \texttt{TankInspector}:\\[2pt]
   \texttt{ValueFarFromTheMedian} (weight)\\
   \texttt{ValueOutsideTheRange} (pH, temperature)\\
   \texttt{ValueMovingTooFast} (pH, temperature)\\
   \texttt{JumpThatComesBack} (conductivity)};

\node[leaf] (powchecks) at (4.3,-4.4)
  {three checks:\\[2pt]
   \texttt{ValueFarFromTheMedian} (current)\\
   \texttt{PairedValuesTooFarApart} (voltage)\\
   \texttt{PairedValuesOutOfProportion}\\(current)};

\foreach \b in {conc, dil, rin} \draw[flow] (\b) -- (tankchecks);
\draw[flow] (pow) -- (powchecks);

\node[wide] (clear) at (0,-7.9)
  {\texttt{Inspector.\_clear\_what\_it\_rejects}\\[2pt]
   \footnotesize \texttt{journal.skip(reading, field, reason)}\quad then\quad
   \texttt{setattr(reading, field, None)}\\[1pt]
   \footnotesize ask the check again; stop at a third of the series};

\draw[flow] (tankchecks) -- (clear);
\draw[flow] (powchecks) -- (clear);

\end{tikzpicture}
\caption{Inside \texttt{SkipSuspectReadings}. The three tank inspectors differ only in which
tank they read; their checks are built once in the shared \texttt{TankInspector}. Every
threshold the checks compare against lives in \texttt{Thresholds}, and a value is cleared in
one place only.}
\label{fig:inspectors}
\end{figure}

\subsection{What each stage is responsible for}

\textbf{\texttt{DataLoader}} is the only part of the program that knows where the data lives.
It reads four tables, validates each against a pandera schema, groups the readings by run,
attaches each lab sample to the tank reading drawn at that minute, and hands back a list of
\texttt{Run} objects. It turns every blank cell into \texttt{None}, because pandas calls a
blank \texttt{NaN} and \texttt{NaN} would quietly survive arithmetic. Which source it uses
comes from \texttt{config.DATA\_SOURCE}, so moving the data out of CSV files is a
configuration change.

\textbf{\texttt{Preprocessor}} decides which runs the modelling may see and which individual
values it may trust. It saves nothing and prints nothing: it returns the surviving runs and
writes what it removed to the journal. It hands on \texttt{Run} objects of exactly the same
type it was given, which is what lets the steps chain.

\textbf{\texttt{Calculator}} is where the readings stop being readings. For each run it
works out the two settings, the two outcomes, and five diagnostics, and packs them into a
frozen \texttt{Result}. Nothing in a \texttt{Result} is copied from a CSV file; every field
was worked out. It drops a run only when one of the two fitted outcomes could not be had,
because a run with no permselectivity and no recovery has nothing to contribute. A missing
diagnostic is not a reason to drop: the command names those runs at the end instead.

\textbf{\texttt{Surrogate}} groups the results by membrane and, for each membrane with at
least three runs, fits all three outcomes and then searches. The fitting is delegated:
\texttt{DesignPoints} logs and codes the settings, \texttt{SecondOrderPolynomial} does the
least squares, \texttt{LeaveOneOutError} scores it, \texttt{ReplicateScatter} works out what
the bench's own repeats cost, and \texttt{FittedOutcome} holds all of it together. It writes
nothing to the journal: a membrane it will not fit is reported on its answer, with the
reason, because a membrane left out names the next experiment worth running.

On the current data this is where equation~\ref{eq:coded} is chosen. \texttt{FittedOutcome}
starts with an intercept alone and adds whichever single term lowers the leave-one-out error
most, stopping when nothing lowers it:

\begin{center}
\begin{tabular}{llrl}
\toprule
model so far & term offered & leave-one-out rmse & \\
\midrule
intercept                 & ---       & 15.0521 & the starting point \\
intercept                 & $x_1$     & 16.2443 & worse \\
intercept                 & $x_2$     & \textbf{7.8979} & best of the two, so $x_2$ is kept \\
intercept, $x_2$          & $x_1$     & 9.9497 & worse \\
intercept, $x_2$          & $x_2^2$   & \textbf{7.1218} & better, so $x_2^2$ is kept \\
intercept, $x_2$, $x_2^2$ & $x_1$     & 9.4710 & worse than 7.1218, so the model stops \\
\bottomrule
\end{tabular}
\end{center}

A squared term may only be offered once its own linear term is already in, since a curve
without its slope cannot be read. That is why $x_1^2$ never appears: $x_1$ never got in.

\textbf{\texttt{BestSetupUnderTheGates}} sweeps a $21 \times 21$ grid of coded points, keeps
those predicted to clear a recovery floor of $35\%$ and a voltage ceiling of
$20\unit{V}$, and returns the survivor with the highest predicted permselectivity. It sweeps
only inside the measured range, so the answer is always a setup the bench has evidence for.
On this data all 441 points fail the recovery floor and none fails the voltage ceiling: the
highest recovery the fit predicts anywhere in range is $33.28\%$. Nothing surviving is
reported as a finding with the blamed gate named, not as an error.

\textbf{\texttt{Analyze}} prints. The exclusions come first on purpose, so a reader who sees
the answer already knows how many runs it came from. Beside each equation it prints the
fit's leave-one-out error next to the bench scatter, and says outright when the error is
worse. For recovery it is: off by $7.122$ against a bench scatter of $6.612$ from two
repeats. The line the command adds is that this is worse than the bench repeating itself,
\emph{so this fit has earned no conclusion}.

\subsection{The one rule: a value that cannot be worked out becomes \texttt{None}}

One rule holds from the loader to the search, and it is the reason the code reads the way it
does. \textbf{A value that cannot be worked out becomes \texttt{None}, and is never guessed
at.} There is no imputation anywhere, no filling forward, no substituting a mean, no
treating a blank as zero.

It shows up at every level:

\begin{itemize}[nosep]
  \item the loader turns a blank cell into \texttt{None}, never $0$ --- a blank means
        \emph{not measured};
  \item a check clears a value it cannot believe, and \texttt{find\_judgeable\_readings}
        then passes over it, so a cleared value is absent rather than wrong;
  \item every quantity returns \texttt{None} rather than a partial answer:
        \texttt{VolumeInATank} with no weight, \texttt{FeedRatio} with no magnesium,
        \texttt{DesignPoints} on a setting at or below zero, which has no logarithm;
  \item \texttt{SecondOrderPolynomial.find\_fit} returns \texttt{None} when the design
        matrix comes back rank-deficient, because two terms moving together over these rows
        means no one set of coefficients is \emph{the} answer;
  \item \texttt{find\_prediction} returns \texttt{None} on an unfitted polynomial, and
        \texttt{BestSetupUnderTheGates} treats a point it cannot predict as a point it
        cannot judge rather than a point that passes.
\end{itemize}

The cost of the rule is that every caller has to check. The benefit is that no number in the
printed report was invented, and a gap in the bench record stays a gap all the way to the
page. The journal is the other half of the same rule: nothing is dropped or cleared without
a line saying what and why.

\subsection{The stages at a glance}

\begin{center}
\footnotesize
\begin{tabular}{@{}>{\raggedright\arraybackslash}p{2.1cm}>{\raggedright\arraybackslash}p{3.7cm}>{\raggedright\arraybackslash}p{3.2cm}>{\raggedright\arraybackslash}p{5.0cm}@{}}
\toprule
stage & the class that does it & takes in & hands on \\
\midrule
load &
  \texttt{DataLoader}, \texttt{CsvSource}, pandera schemas &
  four CSV files &
  21 \texttt{Run} objects, readings and samples attached \\
\addlinespace
preprocess &
  \texttt{Preprocessor} and its four \texttt{Step}s &
  21 \texttt{Run} objects &
  9 \texttt{Run} objects, 10 values cleared; drops and skips in the journal \\
\addlinespace
check a reading &
  \texttt{SkipSuspectReadings}, four \texttt{Inspector}s, six \texttt{Check} classes,
  \texttt{Thresholds} &
  one run's readings &
  the same readings, with unbelievable values set to \texttt{None} \\
\addlinespace
calculate &
  \texttt{Calculator} and the \texttt{quantities} package &
  9 \texttt{Run} objects &
  9 frozen \texttt{Result}s: two settings, two outcomes, five diagnostics \\
\addlinespace
code the settings &
  \texttt{DesignPoints} &
  the settings of each \texttt{Result} &
  one $(x_1, x_2)$ per run, logged, centred and scaled to $-1..1$, and the way back \\
\addlinespace
fit &
  \texttt{FittedOutcome}, \texttt{SecondOrderPolynomial}, \texttt{LeaveOneOutError},
  \texttt{ReplicateScatter} &
  coded rows and one outcome column &
  the terms that earned their place, their coefficients, the error, the bench scatter \\
\addlinespace
search &
  \texttt{BestSetupUnderTheGates} &
  three fitted outcomes and the design &
  the best setup in $\mathrm{mA/cm^2}$ and $\mathrm{mol/mol}$, or the gate that refused
  every point \\
\addlinespace
assemble &
  \texttt{Surrogate}, \texttt{SurrogateModels}, \texttt{Analysis} &
  9 \texttt{Result}s &
  fits and best setups per membrane, plus the membranes not fitted and why \\
\addlinespace
report &
  \texttt{Analyze} &
  the \texttt{Analysis} and the \texttt{Journal} &
  printed lines: exclusions, the run table, the fits, the thin runs \\
\bottomrule
\end{tabular}
\end{center}
